% !TITLE 龟虽寿
% !AUTHOR 曹操
% !DYNASTY 魏晋
% !GENRE 四言古诗
% !ORDER 4

\begin{poem}
神龟虽寿，犹有竟时\textsuperscript{1}；
腾蛇\textsuperscript{2}乘雾，终为土灰。
老骥伏枥\textsuperscript{3}，志在千里；
烈士\textsuperscript{4}暮年，壮心不已。
盈缩\textsuperscript{5}之期，不但在天；
养怡\textsuperscript{6}之福，可得永年。
幸甚至哉，歌以咏志。
\end{poem}

\begin{annotations}
\notetext{1}{竟时：寿命终结的时候。竟，完毕、结束。}
\notetext{2}{腾蛇：古代传说中能腾云驾雾的神蛇。这里用神龟和腾蛇这两种有神力的动物，来比喻任何生命都有终结的一天。}
\notetext{3}{老骥伏枥：骥（jì），能日行千里的好马。枥（lì），马槽。形容有才华的人虽老，依然怀抱远大理想。}
\notetext{4}{烈士：有远大抱负、志向坚贞的人（古代指“立志建功立业的勇士”，非现代烈士之义）。}
\notetext{5}{盈缩：指人寿命的延长与缩短。盈，充实、增长；缩，缩短。}
\notetext{6}{养怡：调养身心，使心情怡然愉快。}
\end{annotations}

\begin{appreciation}
这是曹操《步出夏门行》组诗的第四章，写于北征乌桓胜利之后，是中国文学史上最慷慨激昂的暮年之歌。

曹操写这首诗时已经五十三岁，在古代算是老年了。但他刚刚取得北征的重大胜利，统一北方的大业基本完成，正是雄心勃勃的时候。面对生命的有限和事业的无限，他发出了这首诗中的感慨。

神龟虽寿，犹有竟时；腾蛇乘雾，终为土灰——神龟虽然长寿，也有死亡的一天；腾蛇虽然能乘云驾雾，最终也会化为尘土。曹操首先承认了死亡的必然性。这不是悲观，而是清醒。他不相信神仙长生之说，认为一切生命都有尽头。这种现实主义的态度，在迷信盛行的汉代是非常可贵的。

老骥伏枥，志在千里；烈士暮年，壮心不已——年老的骏马虽然趴在马槽边，心里还想驰骋千里；有抱负的人虽然到了晚年，雄心壮志不会停止。这是全诗最响亮的名句，也是中国文化中不服老精神的经典表达。骥（jì）是千里马，枥（lì）是马槽。一匹老马，即使身体衰弱了，心还在远方。曹操以老马自比，表达了他对生命的热爱和对事业的执着。

盈缩之期，不但在天；养怡之福，可得永年——生命的长短，不完全由天决定；保持身心健康，就可以延年益寿。这里曹操提出了一个积极的人生态度：虽然死亡不可避免，但人可以通过自己的努力来延长生命、提高生命的质量。养怡不只是保养身体，更是修养心性。这种主动掌握命运的信念，体现了曹操作为政治家的务实精神。

这首诗的基调是昂扬的。面对死亡，曹操没有选择消沉或逃避，而是更加珍惜时间、更加努力地追求理想。壮心不已四个字，成为后世无数老年人的精神座右铭。它告诉我们：年龄不是停止追求的理由，只要心还在跳动，梦想就值得继续。
\end{appreciation}
